\documentclass{beamer}

\usetheme{metropolis}

% NOTE: deliberately NOT using booktabs -- all tables are plain tabular + \hline
\usepackage{xcolor}
\usepackage{array}

\hypersetup{
    colorlinks=true,
    urlcolor=blue
}

\title{Product-Key Memory as a Modular Unlearning Module}
\subtitle{Spam unlearning for generative recommendation}
\date{\today}

% ---------------------------------------------------------------------------
% Macros (extend here)
% ---------------------------------------------------------------------------
\newcommand{\pierre}[1]{\textcolor{cyan}{#1}}
\newcommand{\todo}[1]{\textcolor{red}{\textbf{TODO:} #1}}
\newcommand{\done}[1]{\textcolor{green!50!black}{#1}}
\newcommand{\open}[1]{\textcolor{orange!80!black}{#1}}

\newcommand{\If}{\ensuremath{I_t}}      % target (spam) item set
\newcommand{\Df}{\ensuremath{D_f}}      % forget set
\newcommand{\Dr}{\ensuremath{D_r}}      % retain set
\newcommand{\lf}{\ensuremath{\lambda_{\text{forget}}}}
\newcommand{\ls}{\ensuremath{\lambda_{\text{sep}}}}
\newcommand{\gf}{\ensuremath{g_f}}
\newcommand{\gr}{\ensuremath{g_r}}
% per-slot variants: \gf already carries a subscript, so \gf_{,i} would be a
% double subscript -- use these instead.
\newcommand{\gfi}{\ensuremath{g_{f,i}}}
\newcommand{\gri}{\ensuremath{g_{r,i}}}

% status markers for the plan tables
\newcommand{\stdone}{\textcolor{green!50!black}{done}}
\newcommand{\strun}{\textcolor{blue}{running}}
\newcommand{\sttodo}{\textcolor{red}{todo}}

\begin{document}

\maketitle

% ===========================================================================
\begin{frame}{TL;DR}
\begin{itemize}
  \item We can remove $\sim$99\% of spam exposure at \emph{clean-level utility}
        by replacing two decoder FFNs with a PKM and rebuilding them from
        retain data only.
  \item The spam is \textbf{decoder-localised} for removal, but an
        \textbf{encoder} PKM prevents the attack from being learned in the
        first place ($30\times$ less pickup).
  \item \alert{The memory collapses when trained jointly}: 32 of 262{,}144 slots. Post-hoc PKMs (trained in isolation) do not collapse.
  \item Top-$t$ slot selection is implemented but only testable on the non-collapsed (post-hoc) models.
\end{itemize}
\end{frame}

% ===========================================================================
\section{Setup}

\begin{frame}{Setup: model, attack}
\textbf{Model.} TIGER: small encoder--decoder trained from
scratch on semantic-ID sequences ($d_{\text{model}}=128$, 4 encoder and 4 decoder blocks, $d_{ff}=1024$, RQ semantic IDs $L=4$, codebook width 256).

\textbf{Attack.} Bandwagon: fake users click one target item among popular
fillers. Beauty, $1\%$ poisoning, $n_{\text{target}}=1$, target strategy
\texttt{mid} unless stated.

\end{frame}

\begin{frame}{Setup: metrics, reference points}
    
\textbf{Metrics.}
\begin{itemize}
  \item \textbf{SH@$k$} spam hit rate --- fraction of test users seeing \If{}
  \item \textbf{ASI@$k$}, \textbf{TPM@$k$} --- related exposure measures
  \item \textbf{N@10} nDCG (utility), \textbf{UR} $=$ N@10 / clean N@10
  \item \textbf{removal} $= 1 - \text{SH}_{\text{unlearned}}/\text{SH}_{\text{poison}}$
\end{itemize}

\textbf{Reference points (non-PKM), dataset beauty.}
\begin{center}
\begin{tabular}{|l|c|c|}
\hline
model & SH@10 & N@10 \\
\hline
clean            & 0.00000 & 0.03763 \\
poisoned         & 0.01274 & 0.03619 \\
\hline
\end{tabular}
\end{center}
\end{frame}

% ---------------------------------------------------------------------------
\begin{frame}{Why memory? (Geva et al., 2021)}
FFNs behave like \textbf{key--value memories}:
\begin{itemize}
  \item \textbf{keys} capture input patterns
  \item \textbf{values} encode what to output (a distribution over the vocab)
  \item each FFN output is a composition of values; layers refine step by step
\end{itemize}

\medskip
\textbf{Takeaway.} Forgetting by fine-tuning overwrites internal memory
$\Rightarrow$ we want a \emph{modular stabilizer} instead of more backbone edits.

% \medskip
% Our data supports this in a specific way: spam removal only works when
% we touch \emph{decoder} FFNs --- consistent with values encoding the output
% distribution, and generation being decoder-side.
\end{frame}

% ===========================================================================
\section{Design space}

\begin{frame}{Design space: the knobs}
\begin{center}
\begin{tabular}{|l|l|}
\hline
\textbf{axis} & \textbf{options} \\
\hline
When       & trained-in $\;|\;$ post-hoc (insert into trained model) \\
Mode       & \texttt{replace} FFN $\;|\;$ \texttt{add} alongside FFN \\
Where      & encoder / decoder $\times$ block $\{0,1,2,3\}$, or sets \\
Update scope & all params $\;|\;$ \texttt{pkm\_only} $\;|\;$ values only \\
Selection  & none $\;|\;$ AF $\;|\;$ AF-IHF $\;|\;$ grad $\;|\;$ grad+dot \\
$t$        & $1 \ldots 10^4$ (bounded by \emph{live} slots) \\
Optimizer  & Adam $\;|\;$ SGD \\
Init       & \texttt{normal}$\times s$ $\;|\;$ \texttt{zeros} \\
Retain data & sampled subset $\;|\;$ full retain split \\
\hline
\end{tabular}
\end{center}

\medskip
\pierre{All of these are implemented as config knobs; defaults reproduce the
PKM-free behavior.}
\end{frame}

% ---------------------------------------------------------------------------
\begin{frame}{Two variants of the idea}
\textbf{A. Trained-in (``pretrained-memory'').}
PKM is part of the model during training, so it \emph{learns} spam into memory
slots; unlearning edits those slots.

\medskip
\textbf{B. Post-hoc (``correction module'').}
PKM is inserted into an already-trained model, then adapted on clean data as a
sparse corrective layer against a frozen backbone.

\medskip
If PKM is added after training, it has no target information to remove? $\to$ removal is the forgetting, fine-tuning on the retain set repairs this

\end{frame}

% ===========================================================================
\section{Experiments done}

\begin{frame}{Headline comparison (beauty, bandwagon \texttt{mid})}
\begin{center}
\begin{tabular}{|l|c|c|c|}
\hline
method & SH@10 & removal & UR \\
\hline
poison baseline                & 0.01274 & ---    & 0.962 \\
\texttt{seif}                  & $\sim$0 & $\sim$full & 0.67--0.74 \\
\texttt{unified} ac0 (full model) & 0.00411 & 67\% & 1.013 \\
\hline
trained-in PKM, \texttt{pkm\_only} & 0.00385 & 69\% & 1.004 \\
\textbf{post-hoc d01 + rebuild}    & \textbf{0.00013} & \textbf{99\%} & \textbf{0.995} \\
\quad \emph{control:} fresh FFN, no PKM & 0.00022 & 98\% & 0.995 \\
\hline
\end{tabular}
\end{center}

\medskip
\begin{itemize}
  \item Post-hoc \textbf{dominates} everything: clean-equivalent on \emph{both}
        axes (clean is SH 0.00000 / N@10 0.03763).
  \item Trained-in \texttt{pkm\_only} matches full-model unlearning with the
        \textbf{backbone frozen} --- 0.2\% utility cost vs.\ its own clean model.
\end{itemize}
\end{frame}

% ---------------------------------------------------------------------------
\begin{frame}{Placement, post-hoc: spam is decoder-localised}
Insert PKM at \emph{one} layer, rebuild on retain data:
\begin{center}
\begin{tabular}{|l|c|c|c||l|c|c|c|}
\hline
layer & SH@10 & rem. & UR & layer & SH@10 & rem. & UR \\
\hline
d0 & 0.00139 & 89\% & 0.978 & e0 & 0.01306 & $-3\%$  & 0.978 \\
d1 & 0.00031 & 98\% & 0.958 & e1 & 0.01413 & $-11\%$ & 0.978 \\
d2 & 0.00009 & 99\% & 0.904 & e2 & 0.01288 & $-1\%$  & 0.960 \\
d3 & 0.00054 & 96\% & 0.898 & e3 & 0.01400 & $-10\%$ & 0.944 \\
\hline
\end{tabular}
\end{center}

\begin{itemize}
  \item \textbf{Every encoder layer removes nothing.} All decoder layers remove
        89--99\%.
  \item Deeper decoder $\Rightarrow$ more removal, less utility.
  \item \texttt{d01} (both) beats every single layer: 99\% at UR 0.995
        $\Rightarrow$ optimum is the \emph{first two} decoder FFNs.
  \item \texttt{dall} (4 layers) does \emph{not} converge $\Rightarrow$
        a destruction threshold.
  \item Need comparison with FFN reinit instead of PKM replace
\end{itemize}
\end{frame}

% ---------------------------------------------------------------------------
\begin{frame}{Placement, trained-in: encoder PKM is prophylactic}
PKM present \emph{during} training (poisoned run), spam$\times$ relative to
non-PKM (SH 0.01274):
\begin{center}
\begin{tabular}{|l|c|c|c||l|c|c|c|}
\hline
layer & N@10 & SH@10 & spam$\times$ & layer & N@10 & SH@10 & spam$\times$ \\
\hline
D0 & 0.03452 & 0.00277 & 0.22 & E0 & 0.03511 & 0.00072 & 0.06 \\
D1 & 0.03607 & 0.02111 & \alert{1.66} & E1 & 0.03507 & 0.00036 & \textbf{0.03} \\
D2 & 0.03508 & 0.00648 & 0.51 & E2 & 0.03596 & 0.00778 & 0.61 \\
D3 & 0.03480 & 0.01011 & 0.79 & E3 & 0.03508 & 0.00040 & \textbf{0.03} \\
\hline
\end{tabular}
\end{center}

\medskip
\textbf{Encoder PKM suppresses the attack $\sim 30\times$} (E1/E3) for $\sim$3\%
utility. Exactly the opposite of the post-hoc screen --- and consistent:

\begin{center}
\emph{encoder PKM = prophylactic (prevents learning spam)}\\
\emph{decoder PKM = therapeutic (enables removing spam)}
\end{center}

\pierre{Utility spread (0.954--0.997) is inside seed noise; the spam spread
($55\times$) is not.}
\end{frame}

% ---------------------------------------------------------------------------
\begin{frame}{Encoder placements: suppression does \emph{not} stack}
Trained-in, poisoned, decoder untouched. spam$\times$ relative to non-PKM:
\begin{center}
\begin{tabular}{|l|c|c|c||l|c|c|c|}
\hline
layers & N@10 & util$\times$ & spam$\times$ & layers & N@10 & util$\times$ & spam$\times$ \\
\hline
E0   & 0.03511 & 0.970 & 0.06 & E01   & 0.03511 & 0.970 & \textbf{0.04} \\
E1   & 0.03507 & 0.969 & \textbf{0.03} & E012  & 0.03315 & 0.916 & 0.24 \\
E2   & 0.03596 & 0.994 & 0.61 & E12   & 0.03717 & 1.027 & \alert{1.03} \\
E3   & 0.03508 & 0.969 & \textbf{0.03} & E23   & 0.03435 & 0.949 & \alert{2.29} \\
     &         &       &      & E0123 & 0.03739 & \textbf{1.033} & 0.26 \\
\hline
\end{tabular}
\end{center}

\begin{itemize}
  \item \textbf{No combination beats the best single layer} (E1/E3 at 0.03$\times$;
        best combo E01 at 0.04$\times$). Extra blocks buy nothing.
  \item \alert{Combining protective layers can destroy the protection.} E12 gives
        \emph{zero} suppression despite containing E1 (30$\times$ alone); E23
        \emph{amplifies} the attack (2.29$\times$) despite containing E3.
  \item \pierre{Pattern: E2 is a disruptor, E0 a stabiliser --- every combo with
        E2 but without E0 loses suppression; every combo with E0 keeps it.
        Mechanism unexplained.}
  \item Utility does not track layer count either: E0123 (4 blocks) is best
        (1.033), E012 (3 blocks) worst (0.916).
\end{itemize}

\medskip
\textbf{Encoder-PKM suppression is a fragile, non-additive property of specific
blocks --- not a capacity effect.} One configuration makes the model \emph{more}
vulnerable than baseline.
\end{frame}

% ---------------------------------------------------------------------------
\begin{frame}{\alert{Problem: the memory collapses}}
Slot-access diagnostic, $n_{\text{keys}}=512 \Rightarrow 262{,}144$ slots/memory:
\begin{center}
\begin{tabular}{|l|c|c|c|c|}
\hline
model & slots & touched & Jaccard & forget-excl. \\
\hline
trained-in E23D01 & 262144 & \alert{32} & \alert{1.000} & \alert{0} \\
trained-in D1only & 262144 & \alert{32} & \alert{1.000} & \alert{0} \\
trained-in D2only & 262144 & \alert{32} & \alert{1.000} & \alert{0} \\
post-hoc \texttt{replace} & 262144 & 64052 & 0.386 & 3051 \\
post-hoc \texttt{add}     & 262144 & 71547 & 0.411 & 3194 \\
\hline
\end{tabular}
\end{center}

\begin{itemize}
  \item Jointly-trained PKMs use \textbf{0.012\%} of the memory; every token and
        every head picks the \emph{same} 32 slots (routing is input-independent).
  \item \textbf{Post-hoc PKMs are healthy} --- 59\% coverage, forget/retain route
        differently, $\sim$3000 forget-exclusive slots.
  \item \pierre{Hypothesis: trained in isolation on a frozen backbone, the query
        net \emph{must} learn routing; in joint training the model routes
        \emph{around} the memory.}
  \item Consequence: every trained-in PKM number describes a \emph{low-rank
        layer}, not a large sparse memory.
\end{itemize}
\end{frame}

% ---------------------------------------------------------------------------
\begin{frame}{Top-$t$ slot selection: criteria}
Adapt Sparse Memory Finetuning to unlearning --- score slots, update only top $t$.
\begin{center}
\begin{tabular}{|l|l|}
\hline
criterion & score \\
\hline
AF      & forget access frequency (baseline) \\
AF-IHF  & $\text{AF}(s)\cdot\log\frac{T_r+1}{\text{HF}(s)+1}$, history $=$ retain \\
grad    & $\|\gf\| - \lambda\|\gr\|$ \\
grad+dot & $\hat{g}_f - \lambda\hat{g}_r - \mu\,\widehat{\langle \gf,\gr\rangle}$ \\
\hline
\end{tabular}
\end{center}

\textbf{Why the dot term.} The update moves along $+\gf$, so the first-order
change in the \emph{retain} loss from editing slot $i$ is
$\langle \gfi, \gri \rangle$ --- \emph{not} $\|\gri\|$. It is signed:
negative means forgetting also \emph{helps} retain.

\end{frame}

\begin{frame}{Top-$t$ slot selection: criteria}
\textbf{Verified.} Two slots with identical $\|\gf\|=\|\gr\|=12$ but dot $0$
vs.\ $+144$: ranked 1007 vs.\ 1006 by $\gf-\gr$ (indistinguishable), 1007
vs.\ 1999 (last) by grad+dot.

\medskip
\pierre{Objective is additively separable over slots $\Rightarrow$ exact top-$t$
is plain \texttt{topk}, no greedy needed.}
\end{frame}

% ---------------------------------------------------------------------------
\begin{frame}{Selection: gradient signal separates, access counts do not}
Mean per-slot forget/retain gradient cosine:
\begin{center}
\begin{tabular}{|l|c|l|}
\hline
memory & mean cos & reading \\
\hline
E23D01 enc block 2 & $+0.52$ & strongly aligned (conflict) \\
E23D01 enc block 3 & $+0.53$ & strongly aligned \\
trained-in D1only  & $+0.21$ & aligned \\
post-hoc \texttt{d01} & $-0.03$ & \textbf{near-orthogonal} \\
\hline
\end{tabular}
\end{center}

On the healthy post-hoc memory, at $t=25$:
\begin{center}
\begin{tabular}{|c|c|c|c|}
\hline
$\lambda$ & mean $\|\gf\|$ sel. & mean $\|\gr\|$ sel. & overlap with AF \\
\hline
0.0 & 5.06e--3 & 1.88e--2 & 0.36 \\
1.0 & 3.79e--3 & \textbf{1.52e--3} & \textbf{0.00} \\
\hline
\end{tabular}
\end{center}

$\lambda{=}1$ gives up 25\% of forget signal to cut retain exposure
$12\times$, and selects \emph{entirely different} slots than access frequency.
\end{frame}

% ---------------------------------------------------------------------------
\begin{frame}{Ablations that came out negative (worth reporting)}
\begin{center}
\begin{tabular}{|l|c|c|c|}
\hline
post-hoc \texttt{d01} init & SH@10 & removal & UR \\
\hline
\texttt{normal} $\times 1.0$ (default) & 0.00013 & 99.0\% & \textbf{0.995} \\
\texttt{zeros}                          & 0.00009 & 99.3\% & 0.977 \\
\texttt{normal} $\times 0.01$           & 0.00004 & 99.6\% & 0.976 \\
\texttt{add} + \texttt{zeros}           & 0.00018 & 98.6\% & 0.920 \\
\hline
\end{tabular}
\end{center}

\begin{itemize}
  \item ``Smarter'' (identity-like) init is \textbf{worse}: full random init
        rebuilds best. Init \emph{scale} matters, not routing gradients.
  \item Initialising values from the FFN would \emph{re-import the spam} ---
        the removal comes precisely from discarding those weights.
  \item \textbf{Adam beat SGD} in the \texttt{pkm\_only} sweep --- but that
        updated \emph{all} slots, which is not the regime SMF's argument covers.
        Re-test under selection.
  \item \texttt{add}+\texttt{zeros} \emph{also} removes spam (98.6\%): any
        retain-only adaptation suppresses it. \alert{So we cannot easily build
        ``healthy memory $+$ spam still present''.}
\end{itemize}
\end{frame}

% ---------------------------------------------------------------------------
\begin{frame}{Sequential unlearning: no drift, for either method}
Same \emph{total} budget in every arm (chunks $\times$ steps $=$ 3000), same
forget set --- only the number of sequential requests varies:
\begin{center}
\begin{tabular}{|l|c|c|c|c|}
\hline
mode & chunks & steps/chunk & SH@10 & UR \\
\hline
PKM & 1  & 3000 & 0.00031 & 0.981 \\
PKM & 5  & 600  & 0.00031 & 0.992 \\
PKM & 10 & 300  & 0.00022 & 0.992 \\
FFN & 1  & 3000 & 0.00031 & \textbf{1.005} \\
FFN & 5  & 600  & 0.00022 & 0.992 \\
FFN & 10 & 300  & 0.00027 & \textbf{1.005} \\
\hline
\end{tabular}
\end{center}

\textbf{Nothing degrades.} 10 sequential requests cost the same as 1:
97.5--98.2\% removal, UR 0.981--1.005 everywhere.

\medskip
\alert{The modular-memory hypothesis fails on its \emph{premise}:} there is no
accumulated damage for a sparse memory to protect against. A repeatedly
retrained FFN does not drift either (UR 1.005, \emph{above} clean).

\medskip
\pierre{Practical upside: deletion requests can be processed incrementally at no
cost.}
\end{frame}

% ---------------------------------------------------------------------------
\begin{frame}{Collapse is universal in trained-in PKMs}
Slot diagnostic on all five \emph{encoder-only} trained-in models:
\begin{center}
\begin{tabular}{|l|c|c|c|c|}
\hline
model & memories & slots used & Jaccard & forget-excl. \\
\hline
E01   & 2 & 128 / 128        & 1.000 & 0 \\
E012  & 3 & 128 / 95 / 32    & 1.000 & 0 \\
E12   & 2 & 125 / 32         & 1.000 & 0 \\
E23   & 2 & 32 / 32          & 1.000 & 0 \\
E0123 & 4 & 96 / 32 / 32 / 32 & 1.000 & 0 \\
\hline
\end{tabular}
\end{center}

Utilisation 0.012--0.049\% of 262{,}144 slots. \textbf{The hypothesis that the
encoder might escape collapse is refuted} --- collapse is independent of
encoder/decoder placement, layer count and depth. ($32 = knn$ exactly: those
memories return the \emph{same} 32 slots for every token.)

\medskip
\alert{So the E2-disruptor / E0-stabiliser pattern cannot be a memory-content
effect} --- there is no meaningful memory in any of them. Those layers act as
low-rank layers.

\medskip
\pierre{Only healthy memories are post-hoc $\Rightarrow$ it is not \emph{where}
the memory sits but \emph{how} it was trained: alone, against a frozen backbone.}
\end{frame}

% ---------------------------------------------------------------------------
\begin{frame}{Scorecard: PKM vs.\ a fresh FFN}
\begin{center}
\begin{tabular}{|l|l|}
\hline
regime & winner \\
\hline
Converged, single edit       & tie (saturated) \\
Heavy damage (\texttt{dall}) & \textbf{FFN} (UR 0.964 vs.\ 0.891) \\
Small budget                 & \textbf{PKM} (1.75$\times$ removal, $\sim$2$\times$ faster) \\
Sequential edits             & tie (no drift either way) \\
\hline
\end{tabular}
\end{center}

\medskip
\textbf{PKM's only measured advantage is convergence speed at constrained
budgets.} The post-hoc contribution is therefore the \emph{localised retrain}
itself --- re-initialise the first two decoder FFNs, retrain on retain data,
backbone frozen --- with PKM a speed optimisation on top, not the mechanism.

\medskip
\pierre{Remaining lever for the memory story: ``train the memory in isolation,
then unfreeze'' --- the one configuration empirically known to produce a healthy
memory.}
\end{frame}

% ===========================================================================
\section{Plan}

\begin{frame}{What is done}
\begin{center}
\begin{tabular}{|l|c|}
\hline
item & status \\
\hline
PKM in training / post-hoc insertion            & \stdone \\
Memory-only update (\texttt{pkm\_only})          & \stdone \\
Full retain stream (\texttt{retain\_source})     & \stdone \\
Slot access counters + diagnostics               & \stdone \\
Top-$t$ selection (AF / AF-IHF / grad / grad+dot)& \stdone \\
Optimizer + init knobs                           & \stdone \\
Placement screen (post-hoc, 8 layers)            & \stdone \\
Placement screen (trained-in, 8 layers)          & \stdone \\
Architecture study + 168-run unlearning sweeps   & \stdone \\
Clean-baseline audit (seed noise)                & \stdone \\
\hline
\end{tabular}
\end{center}
\end{frame}

% ---------------------------------------------------------------------------
\begin{frame}{What needs to be done}
\begin{center}
\begin{tabular}{|l|l|c|}
\hline
\# & item & status \\
\hline
1 & \textbf{FFN-reinit control} instead of PKM replace   & \stdone \\
2 & Other unlearning algorithms under \texttt{pkm\_only} & \sttodo \\
3 & Fix memory collapse (routing, not capacity)          & \sttodo \\
4 & Top-$t$ sweep on healthy (post-hoc) memory           & \sttodo \\
5 & AF-IHF vs.\ grad vs.\ access baseline ablation       & \sttodo \\
6 & Adam vs.\ SGD \emph{under} selection                 & \sttodo \\
7 & Why does trained-in plateau at SH $\approx$ 0.004?   & \sttodo \\
8 & Multi-seed (clean baselines \emph{and} runs)         & \sttodo \\
9 & Other strategies $+$ other (Amazon) datasets         & \sttodo \\
\hline
\end{tabular}
\end{center}

\medskip
\alert{(1) is the sanity check the post-hoc story rests on:} re-initialise the
same two decoder FFNs as \emph{ordinary FFNs} and fine-tune them identically.
If that matches PKM, the mechanism is ``reinit-and-retrain-this-layer'' and the
PKM is an implementation detail.

\medskip
\textbf{(3) is blocking (4)--(6).} Candidate causes:
\texttt{query\_batchnorm=false} (padding-unsafe), \texttt{query\_proj} collapsing
to a constant, keys never leaving their init.
\pierre{Cheapest lead: post-hoc PKMs do not collapse --- isolate the memory
during training.}
\end{frame}

% ---------------------------------------------------------------------------
\begin{frame}{Item 1: the FFN-reinit control (does PKM actually help?)}
The post-hoc recipe is: \emph{discard} a trained FFN, put a fresh high-capacity
module there, fine-tune it on retain data only.

\medskip
\textbf{Control:} do exactly the same, but re-initialise the layer as an
\emph{ordinary FFN} --- same two decoder blocks, same discovery code, same
optimiser, same steps, same full retain split.
\begin{center}
\begin{tabular}{|l|c|c|c|}
\hline
arm & SH@10 & removal & UR \\
\hline
poison baseline               & 0.01274 & ---    & 0.962 \\
PKM \texttt{d01} $+$ rebuild  & 0.00013 & 99.0\% & 0.995 \\
fresh FFN \texttt{d01} (control) & 0.00022 & 98.2\% & 0.995 \\
\hline
\end{tabular}
\end{center}

\medskip
\alert{Why this decides the framing:} the removal already appears to come from
the \emph{retain-only rebuild}, not from the memory --- ablation alone removes
only 12\%, and \texttt{add} mode (FFN left intact) still removes 98.6\%.
\alert{Result: tied.} Utility identical to 4 decimals (0.03745 vs 0.03746),
removal 99.0\% vs 98.2\%. A plain FFN at \textbf{1/25th} the parameter budget
(2.6M vs 67M) does it just as well.

\medskip
\textbf{So the headline post-hoc result is a \emph{localised retrain}, not a
memory method} --- and it needs no PKM machinery at all.

\medskip
\pierre{But the test is SATURATED: both arms sit at the clean floor. A tie under
saturation shows this task cannot separate them, not that they are equivalent.
See the next slide.}
\pierre{Implemented as \texttt{model.ffn\_reinit\_layers} $+$
\texttt{unlearning.update\_scope=ffn\_only}; the reinit runs AFTER the ckpt load
(FFN weights are in the ckpt, PKM's are not).}
\end{frame}

% ---------------------------------------------------------------------------
\begin{frame}{Where could post-hoc PKM still beat a fresh FFN?}
The tie was measured in a \textbf{saturated} regime --- both arms reach the clean
floor, so the task has no headroom to rank them. To find out whether the sparse
memory is ever worth its 25$\times$ parameter budget, we need a setting where
\emph{neither} saturates:

\begin{center}
\begin{tabular}{|l|l|}
\hline
regime & why PKM might win \\
\hline
More destruction (\texttt{dall}, 4 layers) & more capacity to rebuild with \\
Smaller step budget & which recovers \emph{faster}? \\
Stronger attack / other strategies & harder target to erase \\
\textbf{Sequential requests} & \textbf{repeated edits without drift} \\
\hline
\end{tabular}
\end{center}

\medskip
\alert{The sequential setting is the real test of the modular claim:} a
repeatedly-retrained FFN should accumulate drift, whereas a sparse memory can in
principle localise each edit. That is the SMF/continual-learning argument, and it
is the one regime we have not touched.
\end{frame}

% ---------------------------------------------------------------------------
\begin{frame}{Answer: the two regimes disagree}
\textbf{(A) More destruction --- \texttt{dall}, 4 decoder layers:}
\begin{center}
\begin{tabular}{|l|c|c|c|}
\hline
arm & SH@10 & removal & UR \\
\hline
PKM \texttt{dall} @2000   & 0.00000 & 100\%  & 0.891 \\
PKM \texttt{dall} @10000  & 0.00000 & 100\%  & 0.686 \\
\textbf{FFN \texttt{dall} @2000}  & 0.00009 & 99.3\% & \textbf{0.964} \\
FFN \texttt{dall} @10000  & 0.00004 & 99.6\% & 0.877 \\
\hline
\end{tabular}
\end{center}
\alert{FFN wins by 7 points} and degrades less with extra steps. Where capacity
should have helped most, it \emph{hurt}.

\medskip
\textbf{(B) Smaller budget --- \texttt{d01}, same destruction:}
\begin{center}
\begin{tabular}{|l|c|c|c|c|}
\hline
steps & PKM SH@10 & PKM UR & FFN SH@10 & FFN UR \\
\hline
200        & \textbf{0.00555} & 0.977 & 0.00863 & 0.966 \\
600        & \textbf{0.00148} & 0.989 & 0.00523 & 0.979 \\
converged  & 0.00013 (1560 st.) & 0.995 & 0.00022 (3320 st.) & 0.995 \\
\hline
\end{tabular}
\end{center}
\alert{PKM removes $\sim$1.75$\times$ more at every budget} and converges in
$<$half the steps.

\medskip
\begin{center}
\textbf{PKM buys SPEED, not final QUALITY.}
\end{center}
\pierre{Budget-constrained $\Rightarrow$ PKM. Converged, or heavy damage
$\Rightarrow$ FFN (simpler, 25$\times$ fewer params). Single seed; the
\texttt{dall} gap wants replication, the budget gap is more robust.}
\end{frame}

% ---------------------------------------------------------------------------
\begin{frame}{Item 7: what is the ``SH $\approx$ 0.004 plateau''?}
On the trained-in E23D01 model \emph{every} gradient-based route stops at the
same residual, regardless of parameter scope or budget:
\begin{center}
\begin{tabular}{|l|c|c|}
\hline
route & SH@10 & UR \\
\hline
\texttt{unified}, full model   & 0.00411 & 1.013 \\
\texttt{pkm\_only}, 40 epochs  & 0.00452 & 1.017 \\
\texttt{pkm\_only}, 100 epochs & 0.00385 & 1.004 \\
\texttt{pkm\_only}, 200 epochs & 0.00496 & 1.011 \\
\hline
\end{tabular}
\end{center}

\textbf{Not a measurement floor:} this model's \emph{clean} SH@10 is 0.00004 ---
the plateau sits $\sim$85$\times$ above it, so real spam remains.
\textbf{Not irreducible:} post-hoc reaches 0.00013 on the same attack.

\medskip
Three candidate explanations, each with a different implication:
\begin{enumerate}
  \item \textbf{Optimisation} --- more/better steps would get there.
        \pierre{Weakened: 40$\to$100$\to$200 epochs is non-monotone.}
  \item \textbf{Redundancy} --- spam spread over parameters the objective cannot
        move simultaneously.
  \item \textbf{It is the utility-preserving frontier}, not a ceiling: pushing SH
        lower costs utility. \pierre{\texttt{seif} reaches SH$\approx$0 at
        UR 0.67--0.74, consistent with this.}
\end{enumerate}

\medskip
\alert{If (3), ``ceiling'' is the wrong word} --- the question becomes why the
post-hoc rebuild escapes the frontier entirely.
\end{frame}

% ---------------------------------------------------------------------------
\begin{frame}{Open design question}
The selection experiment needs a model with
\textbf{(a)} healthy memory \emph{and} \textbf{(b)} spam still present.

\medskip
But every route to (a) currently destroys (b):
\begin{itemize}
  \item \texttt{replace} + rebuild $\Rightarrow$ healthy, spam gone (99\%)
  \item \texttt{add} + warmup $\Rightarrow$ healthy, spam gone (98.6\%)
  \item trained-in $\Rightarrow$ spam present, memory collapsed
\end{itemize}

\medskip
\textbf{Options:}
\begin{enumerate}
  \item Fix collapse in trained-in models (then (a) and (b) coexist naturally)
  \item Two-stage: warm \emph{briefly} (routing forms, spam survives)
        $\rightarrow$ select $\rightarrow$ train selected slots only
  \item Reframe: selection as a way to make the \emph{rebuild} cheaper/safer
        rather than to target spam
\end{enumerate}

\medskip
\todo{decide which framing the paper takes}
\end{frame}

% ---------------------------------------------------------------------------
\begin{frame}{Threats to validity}
\begin{itemize}
  \item \textbf{Single seed everywhere.} Three identical-config clean runs span
        \textbf{7.5\%} (sd 3.6\%) --- larger than most architecture effects.
        Utility claims below $\sim$5\% are not currently supportable.
  \item \textbf{One attack, one dataset, one target strategy}
        (beauty / bandwagon / \texttt{mid}).
  \item \textbf{Denominators.} UR must be taken against each model's \emph{own}
        clean baseline; mixing them conflates architecture cost with unlearning
        cost.
  \item \textbf{One cell has no attack} (ALLv2/popular: poison 0.00018 vs.\
        clean 0.00013) --- must be excluded from removal claims.
\end{itemize}
\end{frame}

% ---------------------------------------------------------------------------
\begin{frame}{Resources}
\begin{itemize}
  \item ERASE: \url{https://arxiv.org/abs/2603.08341}
  \item FFNs are key--value memories: \url{https://aclanthology.org/2021.emnlp-main.446/}
  \item Product-key memory (Lample et al.\ 2019): \url{https://arxiv.org/abs/1907.05242}
  \item Memory layers at scale: \url{https://github.com/facebookresearch/memory}
  \item GRID: \url{https://github.com/snap-research/GRID}
\end{itemize}

\medskip
\textbf{Internal.} \texttt{WORKFLOW.md} sections F--J (full experiment record,
job ids); \texttt{tables/pkm\_update.md} (extended notes).

\medskip
\textbf{Deadline.} ICLR 2027 --- Sep 25, 2026 AOE.
\end{frame}

% ===========================================================================
% Add new frames below this line.
% ===========================================================================

\end{document}
